\documentclass[11pt]{article}
\usepackage[a4paper,margin=27mm]{geometry}
\usepackage{amsmath,amssymb,amsthm,booktabs,tabularx,longtable,array}
\usepackage[T1]{fontenc}
\usepackage{lmodern}
\usepackage[colorlinks=true,linkcolor=blue,urlcolor=blue,citecolor=blue]{hyperref}
\newtheorem{theorem}{Theorem}
\newtheorem{lemma}[theorem]{Lemma}
\newtheorem{corollary}[theorem]{Corollary}
\newcommand{\A}[1]{\textnormal{A#1}}
\newcommand{\N}{\mathbb N}
\newcommand{\Z}{\mathbb Z}
\title{Multiplier classifications resolving Sequence Machine's conjecture\\
in OEIS A322026, revision 41, final Formula (8 September 2024):\\
extensions of halasz\_resonator's inverse calculation (2026)}
\author{Haobo Ma\thanks{ChronoAI Pte Ltd, Singapore.}\and
Wenlin Zhang\thanks{National University of Singapore; ChronoAI Pte Ltd, Singapore.}}
\date{7 November 2026}
\begin{document}
\maketitle
\begin{abstract}
Applying and extending halasz\_resonator's inverse calculation (2026,
Theorem 4(1)), we prove Sequence Machine's conjecture
$\A{322026}(n)=\A{323884}(25n)$,
recorded by Antti Karttunen in the final Formula of OEIS A322026,
revision 41, 8 September 2024. The rank function $a=\A{322026}$ is unchanged
on multiplying its argument by a prime greater than $3$. We show that this
invariance forces its Dirichlet inverse $b=\A{323883}$ to vanish at every
multiple of the square of such a prime. We classify all positive multipliers:
$\A{323884}(Mn)=a(n)$ for every $n\geq1$ exactly when
$\gcd(M,6)=1$ and $M$ is not squarefree. The same argument proves and
extends Karttunen's recorded 2022 identity
$\A{126760}(n)=\A{323882}(4n)$.
We prove the neighboring rank and ordinal-transform identities and
distinguish the smooth-number rank array from the rational rank array
in a cross-referenced entry.
\end{abstract}

\section{The sequence and the conjecture}
The final Formula in \cite{a322026} reads
\begin{quote}
\texttt{a(n) = A323884(25*n). [Conjectured by Sequence Machine]}
\end{quote}
Here $b=\A{323883}$ is the Dirichlet inverse of $a=\A{322026}$, and
$C=\A{323884}=a+b$ is a pointwise sum \cite{inverse,sum}.
The conjecture follows once $a(25n)=a(n)$ and $b(25n)=0$.
The inverse mechanism specializes a prior calculation \cite{tower};
the unrestricted multiplier classifications below are its applications
and extensions, rather than a new general inversion principle.
The latter vanishing is a consequence of Dirichlet inversion, rather than
of multiplicativity of $a$: indeed $a(6)=5\ne a(2)a(3)=6$.

For $n\geq1$ define its $3$-smooth part and its complementary factor by
\[
 s(n)=2^{v_2(n)}3^{v_3(n)},\qquad q(n)=n/s(n),\qquad \gcd(q(n),6)=1.
\]
Let $H_1<H_2<\cdots$ list the positive $3$-smooth numbers $2^i3^j$,
where $i,j\geq0$. The rank formula in the posing entry is
\begin{equation}\label{rank}
 a(n)=R(s(n)),\qquad
 R(x)=\#\{(i,j)\in\N^2:2^i3^j\leq x\}.
\end{equation}
Thus $s=\A{065331}$, $R=\A{071521}$ and $H=\A{003586}$
\cite{smoothpart,count,smooth}.
This also proves that \eqref{rank} is the lexicographically earliest
sequence in the name of \A{322026}. A valuation pair $(i,j)$ first
occurs at the integer $2^i3^j$. In a lexicographically earliest labeling,
the first new pair receives the smallest positive label not already
used, and every later occurrence of that pair receives the same label.
Otherwise changing that label to the smallest permissible one gives a
lexicographically smaller admissible sequence. The first occurrences
are therefore labeled in increasing order $H_1,H_2,\ldots$, exactly as
in \eqref{rank}. In particular,
\begin{equation}\label{invariance}
 a(1)=1,\qquad a(Mn)=a(n)\quad\text{if }\gcd(M,6)=1.
\end{equation}

\section{Dirichlet inversion at an insensitive prime}
For integer-valued functions on the positive integers write
\[
 (f*g)(n)=\sum_{d\mid n}f(d)g(n/d),\qquad
 \varepsilon(n)=\begin{cases}1&n=1,\\0&n>1.\end{cases}
\]
If $f(1)=1$, its inverse $g$ exists uniquely: $g(1)=1$ and, for $n>1$,
\begin{equation}\label{recurrence}
 g(n)=-\sum_{\substack{d\mid n\\d>1}}f(d)g(n/d).
\end{equation}
This recursion gives $f*g=\varepsilon$ by construction. All sums are
finite, so the changes in order used below require no analytic convergence.

Prime-variable formal power series for this convolution are classical
\cite[Section 14]{cashwell}. The fixed-prime calculation in
halasz\_resonator \cite[Section 3.4, Theorem 4(1), first proof paragraph]{tower}
already contains the following lemma: if
$F_j=\sum_{p\nmid m}f(p^jm)m^{-z}$, prime-insensitivity makes
$F_j=F_0$, hence $(\sum_{j\geq0}F_jX^j)^{-1}=(1-X)F_0^{-1}$.
Coefficient comparison gives the sign and vanishing rules. We give an
elementary convolution proof of this specialization.

\begin{lemma}\label{support}
Let $p$ be prime, $f(1)=1$, and $f(pn)=f(n)$ for every $n\geq1$.
For $g=f^{-1}$ and $p\nmid m$,
\[
 g(pm)=-g(m),\qquad g(p^r m)=0\quad(r\geq2).
\]
\end{lemma}
\begin{proof}
Let $\delta_p$ be $1$ at $p$ and $0$ elsewhere, and set
$h=(\varepsilon-\delta_p)*f$. Directly,
\[
 h(n)=f(n)-\mathbf1_{p\mid n}f(n/p).
\]
Hence $h(1)=1$, and $h(n)=0$ whenever $p\mid n$.
Put $k=h^{-1}$. We first show $k(n)=0$ whenever $p\mid n$,
by strong induction on $n$. In \eqref{recurrence} for $h$ and $k$,
each summand with $p\mid d$ has $h(d)=0$; in each remaining summand,
$p\mid n/d$ and $n/d<n$, so $k(n/d)=0$ by induction. This proves
the assertion. Commutativity and associativity of convolution now give
\[
 g=k*(\varepsilon-\delta_p),\qquad
 g(n)=k(n)-\mathbf1_{p\mid n}k(n/p),
\]
since multiplying the first expression by $f$ gives $k*h=\varepsilon$.
For $p\nmid m$, this says $g(m)=k(m)$, $g(pm)=-k(m)$, and
$g(p^r m)=0$ for every $r\geq2$.
\end{proof}

\begin{theorem}\label{main}
For every $n\geq1$,
\[
 \A{322026}(n)=\A{323884}(25n).
\]
More generally, if $b=a^{-1}$ and $n=s t$ with $s$ $3$-smooth and
$\gcd(t,6)=1$, then
\begin{equation}\label{mobius}
 b(st)=b(s)\mu(t),
\end{equation}
where $\mu$ is the M\"obius function.
\end{theorem}
\begin{proof}
By \eqref{invariance}, Lemma~\ref{support} applies to every prime $p>3$.
In particular $b(25n)=0$, and $a(25n)=a(n)$, proving the conjecture.
If $t$ has a squared prime factor, the same lemma gives $b(st)=0$
and $\mu(t)=0$. Otherwise write $t=p_1\cdots p_r$ with distinct primes
$p_j>3$. Repeatedly applying $b(p_j m)=-b(m)$, where $p_j\nmid m$,
gives $b(st)=(-1)^r b(s)=\mu(t)b(s)$. This includes $t=1$.
\end{proof}

\begin{corollary}\label{multipliers}
For any positive integer $M$, the equality
$C(Mn)=a(n)$ holds for all $n\geq1$ if and only if
$\gcd(M,6)=1$ and $M$ is not squarefree.
The smallest such multiplier is $25$.
\end{corollary}
\begin{proof}
If the equality holds, put $n=25$. Lemma~\ref{support} gives
$b(25M)=0$, so $1=a(25)=a(25M)=R(s(M))$.
The only $3$-smooth number with rank $1$ is $1$, hence $s(M)=1$
and $\gcd(M,6)=1$. Conversely, if $\gcd(M,6)=1$ and
$p^2\mid M$ for a prime $p>3$, Lemma~\ref{support} gives
$b(Mn)=0$ for every $n$, and \eqref{invariance} gives $a(Mn)=a(n)$.
If $M$ is squarefree, \eqref{mobius} and $b(1)=1$ give
$C(M)=1+\mu(M)\ne1=a(1)$, so $n=1$ is a counterexample.
A nonsquarefree integer coprime to $6$ is divisible by $p^2$ with
$p\geq5$, and hence is at least $25$; the value $25$ works.
\end{proof}

\section{Neighboring formulas in A322026}
\begin{corollary}\label{neighbors}
All nonconjectural Formulas of A322026 hold. Its record-position and
mutual ordinal-transform statements and all four equality-implication
Comments hold as well.
\end{corollary}
\begin{proof}
The rank description also accounts for the other formulas in the posing
entry. David A. Corneth's formulas of 3--4 December 2018 state
\begin{equation}\label{corneth}
 a(H_r)=a((6k+1)H_r)=a((6k-1)H_r)=r\quad(k\geq1),\qquad
 H_{a(n)}=s(n).
\end{equation}
They follow from $s(H_r)=H_r$, $R(H_r)=r$, and
$\gcd(6k\pm1,6)=1$. The record positions are exactly $H_r$:
the first occurrence of label $r$ is at $H_r$, and these first
occurrences have increasing labels. Thus \eqref{corneth} recovers
the already recorded rank and record statements. The precise criterion is
$a(k)>a(i)$ for all $1\leq i<k$, and it holds exactly when $k=H_r$.
The preceding sentence in the source's record Comment has ambiguous
quantifier scope; its following record-position sentence has this precise meaning.

The other formula in Karttunen's 2024 block is
\begin{equation}\label{iverson}
 a(n)=\sum_{j=1}^n[\A{126760}(j)=\A{126760}(n)].
\end{equation}
To see this, write $u(n)=\A{126760}(n)$. Its defining recurrence
\cite{complement} removes factors $2$ and $3$ and assigns
$u(6k+1)=2k+1$, $u(6k+5)=2k+2$. These values enumerate the
integers coprime to $6$ injectively. Consequently $u(j)=u(n)$
if and only if $q(j)=q(n)$. Such $j\leq n$ are precisely
$j=q(n)2^i3^\ell$ with $2^i3^\ell\leq s(n)$; their number is
\eqref{rank}, proving \eqref{iverson}.

If the ordinal transform of $v$ at $n$ denotes
$\sum_{j\leq n}[v(j)=v(n)]$, then $a$ and $u$ are mutual ordinal
transforms, as recorded in the Comments. Indeed $a(j)=a(n)$ means
$s(j)=s(n)$; such $j\leq n$ have the form $s(n)r$ with
$r\leq q(n)$ coprime to $6$. Their number is the rank of $q(n)$
among integers coprime to $6$, namely $u(n)$. Together with
\eqref{iverson} this proves both directions. The same argument makes
$u$ the ordinal transform of $s$, as stated in A126760.
Karttunen had already explained the mutual-transform relation in the
4 December 2018 discussion of A322026 and recorded it in revision 25
\cite{ordinalhistory}; the argument above supplies an explicit proof.

For completeness, the equality implications in the Comments also follow
from the valuation pairs. Equality of $a$ is equality of
$(v_2,v_3)$, hence implies equality of $\min(v_2,v_3)=\A{122841}$,
$\max(v_2,v_3)=\A{244417}$ and the restricted-growth label of their
unordered pair, $\A{322316}$. Equality of that unordered pair implies
equality of $(v_2+1)(v_3+1)=\A{072078}$ \cite{valuationlinks}.
The formula for $\A{305900}$ makes each non-prime-$>3$ input a separate
class and identifies the primes greater than $3$; all the latter have
valuation pair $(0,0)$. Thus equal $\A{305900}$ labels imply equal $a$
\cite{filter}. This proves the source's existing Comments.
\end{proof}

\begin{corollary}\label{dual}
Let $u=\A{126760}$, $d=u^{-1}=\A{323881}$, and
$D=u+d=\A{323882}$. For $3$-smooth $s$ and $t$ coprime to $6$,
\[
 d(st)=\mu(s)d(t).
\]
For any positive integer $M$, $D(Mn)=u(n)$ for all $n\geq1$
if and only if $M$ is $3$-smooth and $(4\mid M\text{ or }9\mid M)$.
In particular $u(n)=D(4n)=D(9n)$; the first equality is Karttunen's
Formula of 18 April 2022 in A126760 \cite{complement,dualinverse,dualsum}.
Moreover $D(n)>0$ whenever $4\mid n$ or $9\mid n$.
\end{corollary}
\begin{proof}
The recurrence for $u$ gives $u(2n)=u(3n)=u(n)$ and $u(1)=1$.
Apply Lemma~\ref{support} successively at $2$ and $3$ to obtain
$d(st)=\mu(s)d(t)$. If $M$ is $3$-smooth and has a squared prime
factor, $u(Mn)=u(n)$ and $d(Mn)=0$, proving sufficiency.
For necessity, evaluate a putative identity at $n=4$:
$d(4M)=0$ implies $u(M)=u(4M)=u(4)=1$.
The rank interpretation of $u$ gives $q(M)=1$, so $M$ is $3$-smooth.
If it is squarefree, $M\in\{1,2,3,6\}$ and
$D(M)=1+\mu(M)\ne1=u(1)$. This proves the classification.
Finally $D(n)=u(n)\geq1$ whenever $4\mid n$ or $9\mid n$.
\end{proof}

The positivity statement restricts possible zeros, but does not settle the
2021 observation in A323882 that zeros apparently have $\mu(n)=-1$,
nor turn its finite nonnegativity observation into an all-index theorem.
The divisor-sum Formula there follows immediately from
$u*d=\varepsilon$: for $n>1$,
$D(n)=-\sum_{1<e<n,\ e\mid n}u(e)d(n/e)$.

\begin{corollary}\label{elementary}
For positive integers $x,n$,
\begin{align*}
R(x)&=\sum_{i=0}^{\lfloor\log_2x\rfloor}
 \bigl(\lfloor\log_3(x/2^i)\rfloor+1\bigr)\\
&=\sum_{k=1}^x\mu(6k)\lfloor x/k\rfloor
 =\sum_{k=1}^x\left(\left\lfloor\frac{6^k}{k}\right\rfloor
 -\left\lfloor\frac{6^k-1}{k}\right\rfloor\right),
\qquad s(n)=\gcd(n,6^n).
\end{align*}
These are the row-count, Cloitre (2007), Browne (2016), and
Sykora (2016) identities in A071521 and A065331 \cite{count,smoothpart}.
\end{corollary}
\begin{proof}
Counting one lattice row at a time gives the first sum.
The sum $\sum_{k\mid m}\mu(6k)$ is $1$ when $m$ is $3$-smooth
and $0$ otherwise: only squarefree $k$ coprime to $6$ contribute,
and their sum is $\sum_{k\mid q(m)}\mu(k)$.
Sum this identity for $m\leq x$ and interchange the finite sums.
The floor difference is $[k\mid6^k]$, again the indicator of a
$3$-smooth integer. Indeed each valuation of $k$ is at most $k$,
so every $3$-smooth $k$ divides $6^k$; the converse follows from
prime factorization. The same valuation bound proves the gcd identity.
\end{proof}

\begin{corollary}\label{arrays}
Writing $T(i,j)=R(2^j3^i)$ for $i,j\geq0$, downward antidiagonal
reading gives A247714, while the transposed array gives A255975.
The identification of this array with A087465 in that entry's
Comment of 25 January 2018 is false: $T(7,0)=49$, whereas
the A087465 array has entry $48$ at $(7,0)$.
The corrected rank parameter is $\log(3)/\log(2)$, not $3/2$.
\end{corollary}
\begin{proof}
A247714 is the position of the antidiagonally listed $2^j3^i$
in the increasing list of $3$-smooth integers, which is $T(i,j)$
by definition \cite{arrays}. A255975 lists the same ranks in the
transposed orientation. Logarithms preserve order, so $T(i,j)$
is the rank of $j+i\log(3)/\log(2)$ among all such values.
In contrast, the definition of the rank array in Kimberling and Brown
\cite[Section 1]{rankarrays} orders $j+(3/2)i$, breaking ties
by the smaller $i$ for A087465. At $(7,0)$ the smooth count is
$12+10+8+7+5+4+2+1=49$.
The rational rank is
\[
1+\#\{(j,i)\in\N^2:2j+3i<21\}
 +\#\{(j,i)\in\N^2:2j+3i=21,\ i<7\}
=1+44+3=48.
\]
This is the zero-based flattened position $35$, where the published
prefixes of A247714 and A087465 give $49$ and $48$, respectively.
\end{proof}

The assertion corrected here is the smooth-rank identification in the
2018 Comment: ``Then R(n,k) is the rank of T(n,k) when all the numbers
in \{T(n,k)\} are jointly ranked.'' Its preceding display has
$T(i,j)=2^j3^i$. This is not a correction to the rational-rank definition of A087465.
Its rational-array recurrence is outside this note's scope.

\begin{corollary}\label{binary}
Let $v(n)=\A{003602}(n)=(n/2^{v_2(n)}+1)/2$ and
$E=v+v^{-1}=\A{349135}$. Then $E(Mn)=v(n)$ for every positive $n$
exactly when $M=2^r$ with $r\geq2$. In particular $E(4n)=v(n)$,
as recorded by Karttunen on 7 December 2021 \cite{binarysum}.
For $w(n)=1+v_2(n)=\A{001511}(n)$ and
$J=w+w^{-1}=\A{323885}$, the identity $J(Mn)=w(n)$ for every
positive $n$ holds exactly when $M$ is odd and nonsquarefree \cite{rulersum}.
\end{corollary}
\begin{proof}
The function $v$ is insensitive to $2$. Lemma~\ref{support} proves
sufficiency for $M=2^r$, $r\geq2$.
Conversely, $n=4$ forces $v(M)=1$, so $M$ is a power of $2$;
$n=1$ excludes $M=1,2$ because the inverse values there are $1$ and $-1$.
For $w$, every odd prime is insensitive. Evaluating a putative identity
at $n=9$ gives $1+v_2(M)=1$, forcing $M$ odd; $n=1$ excludes squarefree
odd $M$, for which the inverse is $\mu(M)$ by Lemma~\ref{support}.
An odd squared prime factor proves sufficiency by the same lemma.
\end{proof}

The question ``Are all terms nonnegative?'' in A349135 was
already answered negatively by Karttunen (30 November 2024): for example,
$E(72747675)=-14126242$. The exact-integer recurrence calculation in
the accompanying checker reproduces this value. This is a prior answer,
not a new nonnegativity conjecture or its first counterexample.

\section{Source coverage}
Here ``proved'' concerns the quoted conjecture or a stated source identity;
an unposed generalization is identified separately. Existing formulas keep
their original attribution. The table covers the Name, mathematical
Comments and every Formula of A322026, together with the relevant companion
items used above; it is not a claim about every formula in the OEIS link graph.
\small
\begin{longtable}{@{}>{\raggedright\arraybackslash}p{.46\linewidth}>{\raggedright\arraybackslash}p{.49\linewidth}@{}}
\toprule Source item and locator & Coverage \\
\midrule\endhead
A322026, rev. 41, final Formula (8 Sep 2024), ``Conjectured by Sequence Machine'' & Proved: Theorem~\ref{main}. \\
A322026, Name; restricted-growth Comment; rank Formula (8 Sep 2024) & Proved: \eqref{rank} and the first-occurrence argument. \\
A322026, Corneth's Formulas (3--4 Dec 2018); record-position Comment & Proved: Corollary~\ref{neighbors}; explicit record quantifiers supplied. \\
A322026, Iverson Formula (8 Sep 2024); mutual ordinal-transform Comment & Proved: Corollary~\ref{neighbors}, including \eqref{iverson}; prior explanation: Karttunen, 4 Dec 2018, rev. 25. \\
A322026, four equality-implication Comments; A322316, matching implication & Proved: Corollary~\ref{neighbors}. \\
Prime-local inverse lemma & Reproved; prior calculation: halasz\_resonator (2026), Theorem 4(1). \\
All positive multipliers for A323884 & Not posed; complete classification proved in Corollary~\ref{multipliers}. \\
A323883 and A323884, Names and sum Formula & Defining inverse/sum identities used in Theorem~\ref{main}. \\
A126760, recurrence, mutual-transform Comments (2019, 2024) & Proved: Corollary~\ref{neighbors}. \\
A126760, Formula $u(n)=\A{323882}(4n)$ (18 Apr 2022) & Proved with attribution, and extended: Corollary~\ref{dual}. \\
All positive multipliers for A323882; dual inverse support & Not posed; classification and factorization proved in Corollary~\ref{dual}. \\
A323882, divisor-sum Formula (18 Aug 2021) & Proved directly from convolution after Corollary~\ref{dual}. \\
A323882, nonnegativity range and apparent zero restriction (18 Aug 2021) & Not covered in full; positivity for multiples of $4$ or $9$ proved. \\
A071521, row-count, Cloitre and Browne Formulas; A065331, Sykora Formula & Proved with attribution: Corollary~\ref{elementary}. \\
A247714, Name; A255975, smooth-rank interpretation & Proved: Corollary~\ref{arrays}. \\
A087465, smooth-rank Comment (25 Jan 2018) & Refuted-and-corrected: $49\ne48$ at $(7,0)$; use $\log_2 3$ instead of $3/2$. \\
A349135, $4n$ Formula (7 Dec 2021) & Proved with attribution; full multiplier classification in Corollary~\ref{binary}. \\
A349135, nonnegativity question and Comment (30 Nov 2024) & Answered: no; prior answer Karttunen (2024), reproduced by an exact finite calculation. \\
A323885, inverse-sum definition & Not posed; complete multiplier classification in Corollary~\ref{binary}. \\
A087465, rational-array recurrence; other unrelated cross-entry identities & Not covered. \\
\bottomrule
\end{longtable}
\normalsize

\section{Mathematical evidence and supplementary material}
In the \href{https://github.com/the-omega-institute/trureturing/blob/e0da8fcc80226febc4e7ea0836927bc71d158292/D5/S3/Arith/DirichletInverse/OeisA322026DirichletInverseIdentity.lean}{pinned Lean source},
the fully qualified declaration
\begin{quote}\small\ttfamily
D5.S3.Arith.DirichletInverse.\\
OeisA322026DirichletInverseIdentity.result
\end{quote}
proves
$a(n)=a(25n)+b(25n)$ for all natural indices, with $a(0)=b(0)=0$.
Its $a$ is the lattice count in \eqref{rank}, its $b$ is the
integer-valued Dirichlet inverse. Its finite box
$0\leq i,j\leq s(n)$ contains every feasible pair, since
$2^i\geq i+1$ and $3^j\geq j+1$. Its private helper
\texttt{inverse\_prime\_square\_support} proves the vanishing part of
Lemma~\ref{support} for every prime and every unit-headed integer-valued
function insensitive to that prime. The restricted-growth interpretation,
the sign rule, \eqref{mobius}, the unrestricted multiplier classification,
Corollaries~\ref{neighbors}--\ref{binary}, and the coverage corrections
have written proofs above; they are distinct from that declaration's scope.

The accompanying \texttt{verify.py} uses integer arithmetic to compare
the rank and first-occurrence definitions through $50\,000$, construct
the Dirichlet inverse by \eqref{recurrence}, compare the published
sequence prefixes, and check the target equality for $1\leq n\leq2\,000$.
It also checks \eqref{mobius}, the companion inverse factorization and
all four multiplier classifications through the bounds in
\texttt{verification.json}, the published rank-array prefixes and the
$49\ne48$ example, the prior negative A349135 value, and the prime-support lemma on four further
integer-valued examples. These finite checks supplement the all-index
proofs. The source archive contains this script,
its numerical inputs and output, the LaTeX source, and build instructions.

\begin{thebibliography}{10}
\bibitem{a322026} A. Karttunen, OEIS A322026, revision 41,
8 September 2024, Name, Comments and Formula; conjecture attributed to
Sequence Machine. \url{https://oeis.org/A322026}.
The \href{https://github.com/oeis/oeisdata/blob/f151d82acb28766e6034a902ed9fd0632b16f986/seq/A322/A322026.seq}{fixed entry text}
specifies the version used here. See also
\href{https://sequencedb.net/s/A322026}{Sequence Machine's A322026 formulas}.
\bibitem{inverse} A. Karttunen, OEIS A323883, revision 5,
8 February 2019, Name. \url{https://oeis.org/A323883}.
\bibitem{sum} A. Karttunen, OEIS A323884, revision 7,
8 February 2019, Name. \url{https://oeis.org/A323884}.
\bibitem{smoothpart} R. Zumkeller, OEIS A065331, Name and Formula.
\url{https://oeis.org/A065331}.
\bibitem{count} B. Cloitre, OEIS A071521, Name and Comments.
\url{https://oeis.org/A071521}.
\bibitem{smooth} P. Zimmermann, OEIS A003586, Name.
\url{https://oeis.org/A003586}.
\bibitem{complement} N. J. A. Sloane, OEIS A126760;
recurrence supplied by A. Karttunen, 28 January 2015, Name and Formula.
\url{https://oeis.org/A126760}.
\bibitem{valuationlinks} OEIS A122841, A244417, A322316 and A072078,
Names and Formulas. \url{https://oeis.org/A122841};
\url{https://oeis.org/A244417}; \url{https://oeis.org/A322316};
\url{https://oeis.org/A072078}.
\bibitem{filter} A. Karttunen, OEIS A305900, Formula.
\url{https://oeis.org/A305900}.
\bibitem{dualinverse} A. Karttunen, OEIS A323881, revision 14,
6 January 2021, Name. \url{https://oeis.org/A323881}.
\bibitem{dualsum} A. Karttunen, OEIS A323882, revision 15,
18 August 2021, Comments and Formula. \url{https://oeis.org/A323882}.
\bibitem{arrays} M. Marcus, OEIS A247714, revision 12,
23 September 2014, Name; C. Kimberling, OEIS A255975, revision 12,
27 March 2015, Comments; C. Kimberling, OEIS A087465, revision 29,
24 August 2024, Name and Comment of 25 January 2018.
\url{https://oeis.org/A247714}; \url{https://oeis.org/A255975};
\url{https://oeis.org/A087465}.
All OEIS entry texts used here are fixed at
\href{https://github.com/oeis/oeisdata/tree/f151d82acb28766e6034a902ed9fd0632b16f986/seq}{the same scientific-data revision}.
\bibitem{rankarrays} C. Kimberling and J. E. Brown,
Partial Complements and Transposable Dispersions,
\emph{Journal of Integer Sequences} \textbf{7} (2004), Article 04.1.6,
Section 1 (definition of the rank array).
\url{https://cs.uwaterloo.ca/journals/JIS/VOL7/Kimberling/kimber67.html}.
\bibitem{cashwell} E. D. Cashwell and C. J. Everett,
The ring of number-theoretic functions,
\emph{Pacific Journal of Mathematics} \textbf{9} (1959), 975--985,
Sections 3 and 14.
\url{https://msp.org/pjm/1959/9-4/pjm-v9-n4-p02-p.pdf}.
\bibitem{tower} halasz\_resonator,
Dirichlet-Inverse Extensions of the Vasyunin-Balazard Biorthogonal System:
Head-Tail Matrix-Loewner Inverse Gram Bounds, Algebraic Reciprocal Obstructions,
and Certified Squarefree-Dilate Separation for the Davenport-Heilbronn Function,
\emph{rStation}, Paper 7, version 1, 16 October 2026,
Section 3.4, Theorem 4(1), first proof paragraph.
\url{https://rstation.ai/problems/riemann_hypothesis/papers/7/v1}.
\bibitem{binarysum} OEIS A003602, Name; A. Karttunen, OEIS A349134, Name;
OEIS A349135, revision 30, Formula of 7 December 2021 and
Comment of 30 November 2024.
\url{https://oeis.org/A003602}; \url{https://oeis.org/A349134};
\url{https://oeis.org/A349135}.
\bibitem{rulersum} OEIS A001511, Name; A092673, Name and Formula;
A323885, Name and Formula.
\url{https://oeis.org/A001511}; \url{https://oeis.org/A092673};
\url{https://oeis.org/A323885}.
\bibitem{ordinalhistory} A. Karttunen, OEIS A322026 edit history,
4 December 2018, discussion at revision 24 and revision 25
(12:36:27 EST), mutual ordinal-transform explanation.
\url{https://oeis.org/history?seq=A322026&start=10}.
\end{thebibliography}
\noindent\small License: Creative Commons Attribution 4.0 International
(CC BY 4.0). \url{https://creativecommons.org/licenses/by/4.0/}.
\end{document}
